\section{Completed cusp spectra and their integral lattices}\label{sec:cusps}

The integral lattice in a secondary target is determined by the homotopy
of a spectrum. Computing this lattice requires retaining both the order
of completion and the real descent at a cusp. We specify these operations
before computing the full target at prime level, where it has two real
factors.

Let $S_N=\Z[1/N]$. We use the normalized Hill--Lawson compactification
$\overline{\mathcal M}_0(N)$ over $S_N$, with smooth locus
$\mathcal M_0(N)$ and full cusp divisor $D_N$. Write
\[
 E_N=\TMF_0(N)=\Gamma(\mathcal M_0(N),\cO^{\mathrm{top}}).
\]
Here $E_N$ denotes a spectrum; the notation $\mathcal E_p$ below
denotes a modular form.

\subsection{Completion before the rational quotient}
For $S=\Z[1/N]$, define the completed real Tate spectrum by
\begin{equation}\label{eq:ordered-completion}
 R_S(t)=\left(\operatorname*{holim}_r
        (\KO[1/N])[t]/t^r\right)[t^{-1}].
\end{equation}
The truncation is the finite truncated monoid algebra, whose coefficient
module has the $r$ summands $1,t,\ldots,t^{r-1}$. Thus coefficient
localization precedes completion in the parameter, which in turn
precedes inversion of the parameter. Homotopy groups are rationalized
only when forming the quotient~\eqref{eq:relative}.

\begin{lemma}[The completed homotopy groups]\label{lem:completed-coefficients}
In every degree $j$,
\[
 \pi_jR_S(t)=(\pi_j\KO[1/N])((t)).
\]
In particular, with $e_K=\alpha\beta^K$, $|\alpha|=4$, $|\beta|=8$,
and $\beta$ invertible,
\begin{equation}\label{eq:real-lattice}
 \pi_{8K+4}R_S(t)=S((t))e_K,\qquad K\in\Z.
\end{equation}
\end{lemma}
\begin{proof}
In degree $j$, the truncations have homotopy
$\bigoplus_{a=0}^{r-1}(\pi_j\KO[1/N])t^a$. Because their transition maps
are surjective in every degree, the Milnor $\lim^1$ term vanishes.
The inverse limit therefore has the product of the coefficient groups,
one for each nonnegative power of $t$. The subsequent localization,
being a filtered telescope, has homotopy consisting of series with a
finite lower Laurent bound. This proves the assertion without a
connectivity assumption on periodic $\KO$.
\end{proof}

To make the rationalization convention explicit, put
\begin{equation}\label{eq:bounded-numerator}
 \Q_S((t))_{\mathrm{bd}}:=S((t))\otimes_{\Z}\Q
 =\{a\in\Q((t)): Da\in S((t))\text{ for some }D\in\Z\setminus\{0\}\}.
\end{equation}
A single denominator $D$ must clear all coefficients away from the primes
dividing $N$. For a prime $\ell\nmid N$, the series
$\sum_{j\ge0}t^j/\ell^j$ lies in $\Q((t))$ but not in
$\Q_S((t))_{\mathrm{bd}}$. Likewise, for $N>1$, $\sum_{j\ge0}t^j/N^j$ belongs to
$S((t))$ but not to $\Z((t))[1/N]$. The rationality of each
individual coefficient does not remove either distinction.

There is a natural comparison
\begin{equation}\label{eq:old-new-comparison}
 \KO((q))[1/N]\longrightarrow R_{S_N}(q).
\end{equation}
The coefficient maps on the finite truncations induce a map first on
their inverse limits and then on their Laurent localizations. This map
factors through $[1/N]$, since $N$ is invertible in the target.
We use the forward comparison for inherited level-one classes; no
inverse or equivalence of the two completed targets is asserted.

\subsection{The full punctured neighborhood}
At the standard cusp, the Tate curve carries the subgroup $\mu_N$.
With the completion convention just fixed, its real restriction map is
\begin{equation}\label{eq:standard-map}
 \Phi_N^\infty:E_N\longrightarrow R_{S_N}(q).
\end{equation}

\begin{proposition}[The punctured full-cusp map]\label{prop:full-cusp-map}
There is a natural map of commutative ring spectra
\begin{equation}\label{eq:full-cusp-map}
 \Phi_N^\partial:E_N\longrightarrow\cK_N^\partial
 :=\Gamma\!\left(\widehat{\overline{\mathcal M}_0(N)}_{D_N}^{\,\times},
                    \cO^{\mathrm{top}}\right).
\end{equation}
Here punctured completion means the completed Tate charts with their
cusp parameter inverted, followed by spectral descent. The construction
includes all cusp components and their descent data and is natural
under forgetting level structure.
\end{proposition}
\begin{proof}
Hill--Lawson realize the completed Tate charts by inverse limits of
finite monoid truncations over multiplicative $K$-theory
\cite[Section~5.1, Theorem~5.7 and Corollary~5.8]{HL}. On each chart,
we first localize the coefficients over $S_N$, then complete its
parameter, and finally invert that parameter. The resulting maps
commute with the chart morphisms. The smooth/Tate comparison and
arithmetic gluing
\cite[Propositions~5.9 and~5.12, Section~5.3]{HL} supply restriction
from the smooth locus, and taking sections of the descent diagram gives
\eqref{eq:full-cusp-map}. This construction is distinct from the empty
ordinary open formal subscheme obtained by inverting an ideal of
definition.

A level-forgetting morphism extends to the normalized compactification
\cite[Proposition~3.19 and Theorem~6.1]{HL}. Locally, a cusp parameter
pulls back to a unit times a positive power of a cusp parameter, so the
resulting completion filtrations are cofinal. The pulled-back parameter
becomes invertible after puncturing, and the compatible Tate
realizations therefore give the claimed naturality on actual spectra.
\end{proof}

\begin{definition}[Joint all-cusp invariant]\label{def:full-cusp}
For a torsion class $y\in\pi_{8K+3}E_N$ killed by
$(\Phi_N^\partial)_*$, let $b_N^\partial(y)=b_{\Phi_N^\partial}(y)$,
with target
\begin{equation}\label{eq:joint-value}
 \cV_{N,K}^\partial=
 \frac{\pi_{8K+4}\cK_N^\partial\otimes\Q}
 {\im\pi_{8K+4}\cK_N^\partial+
  (\Phi_N^\partial)_*(\pi_{8K+4}E_N\otimes\Q)}.
\end{equation}
\end{definition}
The numerator and both indeterminacies in this expression are homotopy
groups of global sections. In particular, the definition of a section
of a homotopy sheaf does not make it an integral element of this
denominator.

\subsection{The two real factors at an odd prime}
\begin{lemma}[The full prime-level cusp spectrum]\label{lem:prime-cusps}
Let $p$ be an odd prime and $S=\Z[1/p]$. Using the invariant elliptic
differential $d\log(z)$ on the Tate curve, there is an equivalence
\begin{equation}\label{eq:prime-cusps}
 \cK_p^\partial\simeq R_S(q)\times R_S(t).
\end{equation}
The two maps from the completed level-one cusp have parameters
$q\mapsto q$ and $q\mapsto t^p$, respectively. Their elliptic curves
and subgroups are
\begin{equation}\label{eq:tate-branches}
 (\Tate(q),\mu_p),\qquad (\Tate(t^p),\langle t\rangle).
\end{equation}
Thus~\eqref{eq:prime-cusps} identifies actual integral homotopy groups,
including the real lattice at the prime two.
\end{lemma}
\begin{proof}
Put $A=S((q))$. On the punctured Tate chart, the finite etale group
$\Tate(q)[p]$ is an extension of $\Z/p$ by $\mu_p$. Locally, a cyclic
subgroup of order $p$ either equals $\mu_p$ or projects isomorphically
onto $\Z/p$. In the latter case, its unique element over $1$ is
represented by a root $t$ of $t^p=q$. This fppf-local description
represents the full finite etale subgroup scheme, with algebra
\begin{equation}\label{eq:subgroup-algebra}
 B_p=A\times A[t]/(t^p-q)=S((q))\times S((t)).
\end{equation}
The polynomial extension is etale because $pt^{p-1}$ is a unit.
For the second equality, group a Laurent series in $t$ by its exponent
modulo $p$: each of the $p$ resulting series in $q=t^p$ has a finite
lower bound. The algebra has rank $p+1$, so neither cusp branch has
been omitted.

Before puncturing, the normalized branches are $S[[q]]$ and $S[[t]]$.
They are finite and normal over $S[[q]]$, with the fraction algebras
in~\eqref{eq:subgroup-algebra}, and therefore compute the normalization.
Their cusp widths are $1$ and $p$, as in the level compactification of
\cite[Proposition~3.19]{HL}.

The completed level-one Tate stack retains the constant group $C_2$,
which acts by inversion on the elliptic curve
\cite[Proposition~3.15]{HL}. Inversion fixes each cyclic subgroup as a
subgroup and hence fixes $B_p$, although it would not fix a chosen
generator for $\Gamma_1(p)$-structure. This inversion action persists
in characteristic two: in the coordinate $x=z-1$, its formal action
is $x\mapsto-x/(1+x)$, which over $\F_2$ is
$x+x^2+x^3+\cdots$, not $x$.

We next compute spectral sections of these stack components. Set
$R=R_S(q)$ and $U=(\KU[1/p])[[q]][q^{-1}]$, with the same ordered
completion as in~\eqref{eq:ordered-completion}. The etale $R$-algebra
$R_{B_p}$ associated to~\eqref{eq:subgroup-algebra} has
\[
 \pi_*R_{B_p}=\pi_*R\otimes_A B_p.
\]
Etale algebras, including their morphisms, are classified by their
degree-zero etale algebra even for nonconnective ring spectra
\cite[Theorem~2.12, Corollary~2.13 and Example~2.14]{MM}.
By the same grouping of powers, the explicit algebra
$R_S(q)\times R_S(t)$ has these homotopy groups in every degree
and thus realizes $R_{B_p}$.

On the complex Tate chart, the level pullback is the corresponding
etale $U$-algebra, whose inversion action is conjugation on $\KU$
with the subgroup coefficients fixed. The Tate realization and
smooth/Tate comparison cited above identify this full coherent descent
datum with the complex base change of $R_{B_p}$.

The compatibility of this descent with the specified completion uses
the faithful $C_2$-Galois extension $\KO\to\KU$ and the dualizability
of $\KU$ over $\KO$ supplied by Wood's theorem
\cite[Section~5.1, Example~5.6, Theorem~5.7 and Example~6.1]{MM}.
Consequently, tensoring with $\KU$ commutes with the inverse limit
of the truncations and with Laurent localization, giving
\[
 R\otimes_{\KO}\KU\simeq U.
\]
Faithful Galois descent is preserved under base change. Hence
\[
 R\simeq U^{hC_2},\qquad
 R_{B_p}\simeq(U\otimes_RR_{B_p})^{hC_2}.
\]
Taking sections of the two actual Tate components now proves
\eqref{eq:prime-cusps}. The argument neither interchanges arbitrary
homotopy fixed points with a Laurent colimit nor replaces homotopy
fixed points by invariants of homotopy groups.
\end{proof}

Complexification sends the compatible Bott generators to
\begin{equation}\label{eq:complexification}
 c(\alpha)=2u^2,\qquad c(\beta)=u^4,\qquad
 c(e_K)=2u^{4K+2},\quad |u|=2.
\end{equation}
Thus $e_K/2$ is not integral in either real factor of
\eqref{eq:prime-cusps}, although its complexification is integral.
This distinction makes the real spectrum computation essential to
the witness constructed below.

\subsection{A rational modular form in global homotopy}
\begin{lemma}[Rational global source]\label{lem:rational-source}
For every integer $w$, naturally under level restriction,
\begin{equation}\label{eq:rational-source}
 \pi_{2w}E_N\otimes\Q
 \cong M_w^!(\Gamma_0(N);\Q).
\end{equation}
The edge map is Tate expansion with the invariant elliptic differential.
For $w=4K+2$, its coordinate in a real cusp factor is
\begin{equation}\label{eq:half-normalization}
 \frac12\Exp_c(f)e_K.
\end{equation}
\end{lemma}
\begin{proof}
The stack $X=\mathcal M_0(N)$ is noetherian and separated. After $N$
is inverted, its level-forgetting map is representable finite etale,
so its map to the formal-group stack is flat. Its relative inertia
kernel is trivial: an elliptic-curve automorphism identical on the
formal completion at the identity is identical on the curve.
For example, equality on the completed local ring forces equality
on the function field and hence on the smooth proper curve. This
verifies relative tameness even in characteristics where the full
modular stabilizer has noninvertible order.

The global-sections theorem of Mathew--Meier
\cite[Definition~2.28 and Theorem~4.14]{MM} therefore applies to the
even-periodic refinement and shows that global sections commute with
homotopy colimits. In particular, it justifies the comparison
\[
 E_N\otimes\Q\simeq\Gamma(X,\cO^{\mathrm{top}}\otimes\Q).
\]
With this comparison established, we calculate rational descent.

Over $\Q$, choose a full auxiliary level $m\ge3$ divisible by $N$,
retaining all its components. The resulting smooth modular scheme
$Y$ is affine, since a nonempty cusp divisor has been removed from
each proper modular-curve component. The stack $X_{\Q}$ is $[Y/G]$
for a finite group $G$. Because taking invariants is exact over $\Q$,
$H^s(X_{\Q},\omega^w)=0$ for $s>0$ and every integer $w$.
The rational descent spectral sequence therefore has cohomological
dimension zero and converges to the edge
identification~\eqref{eq:rational-source}, including negative weights.
Sections on the smooth stack allow finite poles at the omitted cusps.

The edge map is natural under Tate restriction. In the complex
coordinate, the expansion is $\Exp_c(f)u^w$;
equation~\eqref{eq:complexification} then gives
\eqref{eq:half-normalization}. The source is an actual rational global
homotopy class, not only a rational homotopy-sheaf section.
\end{proof}

It follows that the standard-cusp value group, in the coordinate $e_K$, is
\begin{equation}\label{eq:standard-value}
 \cV_{N,K}^{\infty}\cong
 \frac{\Q_{S_N}((q))_{\mathrm{bd}}}
 {S_N((q))+
  \{\tfrac12\Exp_\infty(f):f\in M_{4K+2}^!(\Gamma_0(N);\Q)\}}.
\end{equation}
The source image lies in the bounded numerator because it is the
image of $\pi_{8K+4}E_N\otimes\Q$ under an actual spectrum map.
Likewise, the rational source in the joint quotient is obtained by
simultaneously restricting one global class, rather than choosing
a class independently at each cusp.

Let $\cV_{1,K}$ denote the quotient~\eqref{eq:relative} for the
inherited level-one map $\Phi_1:\TMF\to\KO((q))$. Its localization
$\cV_{1,K}[1/N]$ uses the localized inherited target, before the
forward comparison~\eqref{eq:old-new-comparison}.
\begin{corollary}[Naturality from level one]\label{cor:level-naturality}
The forgetful map $\iota_N:\TMF\to E_N$ and cusp restriction induce
\[
 \rho_{N,K}^\partial:\cV_{1,K}[1/N]\longrightarrow\cV_{N,K}^\partial,
 \qquad
 b_N^\partial(\iota_Nx)=\rho_{N,K}^\partial(b_1(x))
\]
for every torsion $x\in\pi_{8K+3}\TMF[1/N]$.
\end{corollary}
\begin{proof}
Apply Proposition~\ref{prop:naturality} to the forward comparison
and the actual level-forgetting restriction square of
Proposition~\ref{prop:full-cusp-map}. Since
$\pi_{8K+3}\KO((q))=0$, the level-one primary image is zero;
the imported class therefore lies in the primary kernel at level $N$.
\end{proof}
